% chapters/02_grundlagen.tex

\chapter{Grundlagen}
\label{chap:grundlagen}

\section{Humanoide Robotik}
\label{sec:humanoide-robotik}

Humanoide Roboter zeichnen sich durch eine dem Menschen ähnliche Körperstruktur
aus, die eine intuitive Interaktion mit menschlichen Werkzeugen und Umgebungen
ermöglicht~\cite{placeholder_humanoid}.
Der \textbf{Unitree G1} verfügt über 43~\acp{dof} und wiegt ca.\ \qty{35}{\kg};
die \textbf{DEX3-Hand} ergänzt diesen um 6~\acp{dof} pro Hand für Feinmanipulation
(siehe \cref{tab:hardware}).

\begin{table}[H]
  \centering
  \caption{Hardware-Spezifikationen: Unitree G1 mit DEX3-Hand}
  \label{tab:hardware}
  \begin{tabularx}{\linewidth}{lX}
    \toprule
    \textbf{Eigenschaft}   & \textbf{Wert} \\
    \midrule
    Gesamtgewicht          & \qty{35}{\kg} \\
    Körper-DoF             & 43 \\
    DEX3-Hand-DoF (je)     & 6 \\
    Kamera(s)              & Stereokamera + Head-Cam \\
    Eigenfrequenz Regelung & \qty{500}{\Hz} \\
    \bottomrule
  \end{tabularx}
\end{table}

\section{Vision-Language-Action-Modelle}
\label{sec:vla}

\acp{vla} sind multimodale neuronale Netze, die Kamerabilder, Sprachbefehle und
Roboterzustände gemeinsam verarbeiten und direkt Steuerungsbefehle ausgeben~\cite{openvla2024}.
\textbf{GR00T N1.6} von NVIDIA folgt einer Dual-System-Architektur~\cite{groot2025}:

\begin{description}
  \item[System~1 (reaktiv)] Ein leichtgewichtiger \ac{vit}-basierter
    Aktions-Tokenizer generiert flüssige Bewegungssequenzen mit niedriger Latenz.
  \item[System~2 (kognition)] Ein großes Sprachmodell verarbeitet Sprachbefehle
    und steuert die übergeordnete Planung (\ac{moe}-Architektur).
\end{description}

\begin{figure}[H]
  \centering
  % \includegraphics[width=0.85\textwidth]{figures/groot_architecture}
  \fbox{\parbox{0.85\textwidth}{\centering\vspace{2cm}
    [Platzhalter: GR00T N1.6 Architektur-Diagramm]\vspace{2cm}}}
  \caption{Dual-System-Architektur von GR00T N1.6 (nach \cite{groot2025}).
    System~1 generiert Aktionen reaktiv; System~2 plant auf semantischer Ebene.}
  \label{fig:groot-architecture}
\end{figure}

\section{Imitation Learning}
\label{sec:imitation-learning}

\ac{il} ermöglicht das Erlernen von Verhaltensweisen aus menschlichen
Demonstrationen ohne explizite Belohnungsfunktion~\cite{hussein2017imitation}.
Das Training minimiert die \emph{Behavior-Cloning}-Verlustfunktion:

\begin{equation}
  \mathcal{L}_{\text{BC}}(\theta)
  = \mathbb{E}_{(s,a) \sim \mathcal{D}}
    \bigl[\lVert \pi_\theta(s) - a \rVert^2\bigr],
  \label{eq:bc-loss}
\end{equation}

wobei $\pi_\theta$ die Modellpolitik, $s$ den Zustand, $a$ die demonstrierte
Aktion und $\mathcal{D}$ den Demonstrations-Datensatz bezeichnen.

\begin{warningbox}[Compounding-Errors-Problem]
  Bei reinem Behavior Cloning können sich kleine Fehler über Zeithorizonte
  aufaddieren (\emph{compounding errors}), da das Modell keine Fähigkeit zur
  Selbstkorrektur erlernt. Datenaugmentation und diverse Demonstrationen
  können diesen Effekt abmildern.
\end{warningbox}

\section{Datensatzformat: LeRobot}
\label{sec:lerobot}

Der Demonstrations-Datensatz liegt im \textbf{LeRobot}-Format (v2.1) vor,
einem von Hugging Face entwickelten Standard für Robot-Learning-Datensätze~\cite{lerobot2024}.
Ein Datensatz besteht aus Episoden mit synchronisierten Kameraframes
(RGB, \qty{30}{\fps}), Gelenkpositionen, Geschwindigkeiten und Aktionsvektoren.

\begin{tipbox}
  Das Format enthält eine \texttt{modality.json}, die alle Sensormodalitäten und
  deren Dimensionen beschreibt. GR00T erwartet v2.1; der im Repository enthaltene
  Konverter übersetzt v3.0~$\to$~v2.1 automatisch.
\end{tipbox}
